\chapter{Perception et navigation autonome}
\label{chap:perception}

\section{Contexte}
% TODO : après la téléopération, rendre le véhicule capable de « voir » et de
% décider seul — et fournir la perception qui alimentera le V2V (chap. 7).

\section{Problématique}
% TODO : percevoir (lidar + caméra), se localiser, naviguer et freiner seul,
% avec les ressources limitées d'un calculateur embarqué.

\section{Mission 1 : Cartographie et localisation (SLAM)}
\objectif
\methodes
% TODO : RPLIDAR -> slam_toolbox, réglages, cartes obtenues.
\resultats
\discussion

\section{Mission 2 : Navigation réactive follow the gap et avancement Nav2}
\objectif
\methodes
% TODO : principe follow_the_gap sur le scan lidar, package navigation,
% intégration twist_mux (priorités avec la téléop) ; état d'avancement Nav2.
\resultats
\discussion

\section{Mission 3 : Détection de piétons par vision --- YOLO embarqué sur GPU}
\objectif
\methodes
% TODO : portage de l'inférence sur le GPU de la Jetson : JetPack/CUDA,
% installation de torch spécifique Jetson, numpy compatible — le chemin de
% croix technique mérite d'être raconté.
\resultats
% TODO : LE tableau avant/après : 344 ms/image (CPU) -> 30 ms/image (GPU),
% soit un facteur ~11 — ton meilleur résultat chiffré du chapitre.
\discussion

\section{Mission 4 : Freinage d'urgence automatique (AEB) par fusion lidar + caméra}
\objectif
\methodes
% TODO : fusion aeb_fusion (offset 180° du lidar, seuils dynamiques selon la
% vitesse, frein actif calibré, timeout twist_mux), « boîte noire » de debug.
\resultats
% TODO : scénarios d'essai (obstacle statique, piéton), distances d'arrêt.
\discussion

\section{Conclusion}
